I want to pursue NTU's MSc in Robotics and Intelligent Systems because real-robot work showed me that a learned policy is only one part of autonomy. While deploying an OpenPI loop on a Dobot CR5, I assembled camera and robot-state observations, then translated returned action chunks into arm and gripper commands. Under non-blocking execution, a new inference could use an intermediate in-motion state before the preceding chunk finished, producing back-and-forth corrections. I chose blocking, short-horizon ServoP execution instead, accepting lower apparent frequency to preserve state-action consistency. This made me want to study how perception, control, planning, and execution should be designed together.

The programme would extend this experience across manipulation and mobile autonomy, letting me compare how sensing and execution constraints change between these two settings. MA6212 Robotics Manipulation and Advanced Control would give me stronger tools for analysing manipulator dynamics, trajectory planning, and control beyond my empirical choice. MA6214 Autonomous Mobile Robot would broaden the comparison to navigation, motion control, localisation, and multimodal fusion. In the coursework-only route, MA6215 Robotics Projects could support a focused prototype examining how observation timing and action-completion rules affect a robot pipeline.

I would bring experience integrating synchronized robot data, policy serving, and physical execution, together with a habit of tracing behaviour to the layer that produced it. My immediate goal is robotics or embodied-AI R\&D, connecting learning to sensing, control, and dependable deployment. NTU would help me move from making one real-robot loop work to reasoning rigorously about autonomous systems across embodiments. After stronger real-system experience, I may consider doctoral study, but my priority is to design and evaluate complete robotic systems.
